\documentclass[12pt,stu,biblatex]{apa7}
\usepackage{parskip}
\setlength{\parindent}{1.5em}
\usepackage{indentfirst}
\usepackage{geometry}
\usepackage{setspace}
% \geometry{left=20mm, top=20mm, bottom=20mm}

\usepackage{hyperref}
\usepackage[backend=biber,style=apa]{biblatex}
\addbibresource{refs.bib}

\author{Avi Chetlin}
\affiliation{Case Western Reserve University}
\course{ORBH 375: Building Leadership Character}
\title{Class Journal, Weeks 1-7}
\duedate{08 October 2026}
\professor{Dr. Susan Case}

\begin{document}
\maketitle
\singlespacing
\tableofcontents
\newpage

\section{Week 1. Leadership Character in Context: Exploring Virtue \textit{(27 August 2026)}}

\subsection{Reading Summary}

In \textit{Defining the Good, the Bad, and the Evil} \parencite{wood2021goodbadevil}, Wood, Meister, and Liu begin by arguing that discussions of bad leadership become clouded when the ideas of effectiveness and morality are treated as one and the same.
When we call a leader ``good,'' we might mean that they accomplished a final goal, or we might mean that the person acts humanely, responsibly, and informed by an ethical code with which we agree.
These are fundamentally different judgments.
A leader may successfully head a team while actively harming others; a person with admirable intentions can lose control of their team or fail to reach their desired outcome.
Wood, Meister, and Liu therefore define the ideas of good and bad leadership entirely focused on the tangibles, that is, good leadership accomplishes the goal at hand, and bad leadership doesn't.
The moral standing of leaders must be assessed separately.
The authors' terminology is intended to prevent this definition of leadership from excluding leaders who are both extremely effective and extremely immoral (pp. 47-51).

The authors also make a distinction between morality and ethics. The former involves judgment based on conscience and a considered sense of right and wrong; the latter concerns the conformity of conduct to common external standards.
This challenges the assumption that accepted professional behavior, institutional permission, or loyalty to a superior necessarily establishes moral responsibility.
It additionally extends responsibility to followers, in the sense that the way in which a person responds to a command reveals something about their own participation in the leadership process.
This reminded me of the concrete example of the soldiers of the SS during the Holocaust.
A common refrain of theirs was that they were, ``just following orders.''
The ethical/moral distinction drawn out by Wood, Meister, and Liu helps explain why following an order does not remove responsibility for the action.
This distinction does not make every difficult case easy, but it prevents obedience from being a complete defense.
This was my favorite part of the reading, and something I'm most likely to apply to my own leadership style; distinguishing between doing the correct or normal thing as laid out by others, the effective thing as determined by economic pressures, and the moral decision in terms of impact on all people who could be involved.

There is also something useful in their treatment of leadership as a process involving other people, rather than a property someone acquires with a title.
A leader needs followers, and followers make choices about which purposes to support and which conduct to accept.
This means that responsibility is distributed, although certainly not equally.
The person with authority can shape the available choices in a way that a subordinate often cannot.
To me, this is a more realistic way to evaluate leadership than treating the person at the top as the sole cause of everything an organization does.
It also leaves room for the kind of contribution I tend to prefer, where I can influence the character of a group without necessarily wanting to be in charge of it.

Newman makes a related distinction between competence, commitment, and character \parencite[ch. 1]{Newman2019building}.
Someone can be intelligent and hardworking while using both qualities toward an objectionable end.
The important question is therefore not just whether someone has the capacity to lead, but what they consider worth doing with that capacity.
I find this helpful in thinking about my own tendency to admire intelligence.
Intelligence is something I value tremendously, but it cannot, by itself, establish whether I should trust a person's judgment about how to treat someone else.

\subsection{Journal Prompt and Discussion}

In all honesty, I had a tough time coming up with a specific person to whom I look up to for their strong character.
I sat with the question for a while, and felt weird about how long it took me to come up with an answer.
In class, it seemed like a lot of people had automatic answers; Oliver, for example, had a very clear understanding of who he looked up to and why, which made me wonder why I didn't have that.
Certainly, there are plenty of people I aspire to resemble in terms of their knowledge, or accomplishments, or kindness, or any number of individual traits, but I think I subconsciously cast judgment on others for their character or integrity more than anything else.
I used to particularly admire my father's character, and I still do, but some recent discussions we've had and some unfortunate circumstances he's encountered have made it more difficult for me to cast him as a pure moral example.

I think part of my difficulty is I'm searching for moral absolutes, a perfect person, which obviously doesn't exist.
My father resembles the kind of man I'd like to be in his integrity, in his empathy, in his commitment to his values, and in his intelligence.
But he's made many mistakes I hope to avoid and hurt people in ways I'd like to prevent myself from doing.
His character, though, is a nice reflection of \parencite{wood2021goodbadevil}; he has tremendous moral character and the capacity to be a strong leader, but hasn't succeeded as a professional leader in terms of his explicit accomplishments yet.

The question of what I'd like to have included in my obituary, however, was easier for me to answer.
I've actually spent a good amount of time thinking about this question myself, and the related question of whether or not it matters the legacy I leave behind.
In fact, I find, in some sense, comfort in the idea that I won't leave very much behind.
Certainly it's true that I'd like to be remembered as someone with kindness, with warmth, with a willingness to help and a continued empathy for others, but I have no real desire to be remembered for massive accomplishments.
Well, sometimes I think that might be from an awareness that I'm in fact nowhere near the raw talent or intelligence of those who will have an opportunity to attach their name to discoveries or big ideas, but let's suppose for the moment I truly don't care about leaving behind a tangible legacy.
Then what I really do care about is how I'm perceived as a person, how I leave my mark, as a human, on others, and how I create a positive impact for those surrounding me.
All of this is to say, I would want my obituary to be written by someone who loved me, saying, ``we all know about Avi's generosity, his empathy, his understanding, but what you might not know is...''

Following the experience and thinking I described at the beginning of this subsection, I spent some time asking others if they had an answer for their character exemplars, and discovered that actually, many of my friends and colleagues don't have specific examples of individual people.
They mentioned that in fact they based their character on an amalgamation of others, on particular moments they remember being influenced by, which sounds much closer to what I'm used to.

The conversations with my friends gave me a small test of the assumption I brought into class.
I had treated the absence of a single exemplar as a deficiency, as though everyone else had found an obvious answer that I was missing.
Asking them about it showed me that the assumption was too strong.
I can learn from a particular act of kindness or integrity without having to defend everything about the person who performed it.
That feels like a more useful starting point for developing my own character, and a fairer way to look at my father.

To conclude this week's entry, I'd like to mention that I have trouble imagining what the most important leadership qualities for me to improve upon are, because I don't particularly want to be a leader.
I have moments of strong leadership, and in fact often find myself in those roles on the smaller scale (class projects, reading groups, organizing plans with groups of friends) and on the medium scale (school organizations, research) but in an ideal world, as put by Aaron Sorkin, I'd rather not be the Guy, I'd rather be the guy that the Guy can depend on.

\section{Week 2. Learning Character Lessons from Failure \textit{(1 September 2026)}}

\subsection{Concrete Experience}
In high school, I did extremely poorly academically, which I was sad and ashamed about.
That experience informed how seriously I take my academics and work now.
Some of it was to do with the pandemic, which tremendously impacted my motivation in school and cratered my social intelligence.
But some of it was just a lack of maturity and understanding of what was important at the point in my life, and that is what I deeply regret.
I wish I could go back in time and prevent myself from making these mistakes, but I can't, so I have to live with this failure.

\subsection{Reflective Observation}
Failure is painful and scary.
In general, failure is to be avoided at most costs.
Failure is disappointing and embarrassing.
Most of the pressure of failure comes from those involved; failing on a solo project mostly makes me disappointed in myself, and rarely is the external pressure anywhere as strong as that I put on myself.
Because I avoid failure, I can, and often do, miss out on opportunities.
However, other people also gain an impression of me as being successful because they aren't familiar with my failures.

I'm obviously invested in my career, academic or otherwise, which warrants some focus on grades and academic performance, but there's a further obsession with perfection that I'll be the first to admit is irrational.
That being said, it hasn't caused me too much undue stress yet, and striving for a 4.0 isn't the worst irrational compulsion of all time, so I don't plan on stopping that yet.

\subsection{Abstract Conceptualization and Reading Reflection}
In the first chapter of \textit{Better Humans, Better Performance}, Rea, Stoller, and Kolp argue that character and competence need to develop together \parencite[ch. 1]{Rea2022better}.
Being capable matters, but capability alone does not determine what a person will do under pressure.
The chapter's examples place considerable emphasis on relationships, purpose, and the conditions in which people are trying to improve.
The EDWINS example is useful in this respect because it combines expectations of performance with opportunities and support.
It would be easy to tell someone that they should respond better to adversity; it is more demanding to create circumstances in which they have a realistic opportunity to do so.

This makes failure a more complicated subject than simply saying that making mistakes builds character.
A failure can expose a missing skill, an unrealistic expectation, or a habit of avoiding an uncomfortable truth, and those are not interchangeable problems.
If the problem is a missing skill, the appropriate response might be further practice or asking for help.
If the problem is hiding a mistake, then becoming more technically capable will not necessarily address what went wrong.
There is an important distinction between the unsuccessful outcome and the person's subsequent response to it.
The latter is where questions of responsibility, honesty, and care for other people become particularly apparent.

In Chapter 2, the authors describe trust, compassion, courage, justice, wisdom, temperance, and hope as virtues that can be developed through repeated practice \parencite[ch. 2]{Rea2022better}.
The value of this framework is that it makes improvement more specific.
Courage might mean admitting a mistake before someone discovers it; justice requires considering who has absorbed its consequences; temperance can make room between an initial emotional response and the action that follows.
Wisdom is needed to determine which response fits the circumstances, since neither persisting nor giving up is automatically the better choice.
Hope, similarly, should make another attempt possible without requiring someone to ignore evidence that the approach needs changing.

Read together, these chapters suggest that learning from failure requires more than identifying a lesson afterward.
The lesson needs to become a change in how the next situation is handled.
At the same time, the fact that a mistake taught someone something does not cancel the harm it may have caused someone else.
Repair and personal growth can both be necessary.
This connects with Newman's distinction between competence, commitment, and character.
An intense desire to succeed is not sufficient if protecting the appearance of success takes priority over acknowledging what happened \parencite[ch. 1]{Newman2019building}.

As I explained earlier, I have a really tough time getting behind failure on principle.
I grew up playing hockey, and we were taught that if we weren't falling down, we weren't trying hard enough.
I also grew up doing musical theatre, where a similar lesson of full effort and failing hard was instilled upon us.
That being said, it's just really tough to intentionally put oneself into a situation wherein failure is possible if not likely, and that's something I'm hoping to work on.

\subsection{Active Experimentation}
This semester I took on a fairly substantial work load, including 19 credits, 9 of which are graduate level, a research project, two jobs, and several club commitments.
I'm excited about each of those components individually, so I decided to do them all instead of playing it safe, but that certainly poses a risk to my perfectionism.
I feel good, so far.
It feels good to be applying my abilities and intelligence, my improvements in executive functioning and overall academic confidence.
I would be lying, though, if I said I wasn't under more pressure than I have been in previous semesters, or that I don't occasionally feel overwhelmed.
There are days I wish I had an easier schedule, but there are days when I feel proud of what I'm accomplishing.

\subsection{Closing Ethical Reflection}
I'm still not sure if I made the right choice or not in choosing this schedule, but so far it seems advantageous and it hasn't caused me to crash and burn yet.

\section{Week 3. Ethical Heroes: Turning Admiration into Practice \textit{(10 September 2026)}}

\subsection{Concrete Experience}
Two summers ago, during my family's vacation in North Carolina, my father told me that he had lost his job early in President Trump's first term.
That is, he told me years after it had happened.
I was, at that point, somewhat aware that something had happened, and in fact that other family members were aware that I hadn't been told and were disappointed in my father as such.
He was deeply ashamed, horrified really, to have to tell me that.
He was embarrassed, not just to have to tell me he'd lost his job but that he had to fess up to keeping it from me for so long.
I remember feeling shocked and, if I'm honest, scared.
I was sad he hadn't told me until then, and I didn't understand how my father, who I think of as one of the smartest people I knew in the world, could be out of a job.

This opening incident is useful to my reflection because it's an ethical role model of mine making a sudden departure from my idealized perception of them to a real-life conflict.
The principles I see at stake are integrity, trust, and courage.
The immediate effect on me was a loss in confidence, and a feeling that maybe my father hadn't trusted me enough to tell me, but also the relief that that secret was no longer underpinning our relationship.
\subsection{Reflective Observation}
Admiration, love, and respect, can make criticism difficult.
Even though I felt somewhat betrayed by my father's admission, scared about the financial implications, and confused about what that meant about the man I thought I knew and looked up to, I found myself immediately trying to explain away what he had told me, to forgive him subconsciously without processing the situation consciously.
Clearly, my father accepted the cost of embarrassment and shame in telling me, while I bore the cost of having to reconcile the image I had of my father with the truth.

To me, the ``right behavior'' is averaging the human cost of your decisions, perhaps lightly weighted by the relative relational importance of the people, and choosing that which harms the fewest.
The virtues I'd like to develop from this example are trust and temperance.
The former is clear; I don't want to keep things from people that might hurt them if they emerge later.
To temperance, though, I'd like to work on controlling my knee-jerk reactions.
As I detailed in the CE section, it's clear some of my immediate reactions could have been harmful to my father, and even though I didn't voice them, I was ashamed of myself for having them.

\subsection{Abstract Conceptualization and Reading Reflection}
John Frye's primer on ethical decision making \parencite{fryeEthicalAnalysis} provides different tests for evaluating the same act.
Consequentialism asks about the effects on those who are affected; deontology asks about duties and respect for persons; virtue ethics asks what a choice expresses about character and practical judgment.
These ideas reminded me a bit of the reading from the first week, \parencite{wood2021goodbadevil}, in that it explored juxtapositions of how different ways to evaluate behavior can explain complicated behavior.
A beneficial outcome can coexist with a broken promise, and a principled intention might coexist with avoidable harm, just like an efficient leader acting immorally, or vice versa.

Care ethics, then, adds the relationship and the recipient's perspective to our evaluation.
Engster and Hamington \parencite{engsterHamington2015introduction} lay out this approach---attentiveness to interdependence, particular needs, responsibility, and responsiveness.
That is, good intentions on the part of the carer do not mean that the care was adequate.
The person receiving help must have room to respond, to speak to whether their needs were met.
This is different from requiring them to return the favor; the point is that the recipient can correct the carer's understanding of what would actually help.

Applied to my example, the conceptual question is whether the behavior of my father had both a defensible purpose and an attention to those affected.
To the former, my father obviously had to tell me eventually, and I'm glad he did.
I also think it's undeniable that anyone involved would agree on that point.
To the latter, though, although I think my father told me out of love and care, and although he expressed a worry about me for having to hear about his situation, I do think he was more focused on himself in that situation.
I think I would be, too, but I don't think he exhibited the ideas from care ethics in this particular case.
This is the evidence on which the ethical judgment should rest.
Where the evidence is incomplete, the conclusion should remain limited to the action observed rather than declaring the whole person virtuous (or not).
I really do think my father is an ethical exemplar in many ways, and this particular situation does not entirely color his character.

The frameworks also help me qualify my own description of right behavior as minimizing human cost.
That description is largely consequentialist, but the other readings raise questions that cannot be answered simply by adding up distress.
Someone might be less upset in the short term because they have not been told something important, while still having a legitimate claim to know it.
Similarly, giving additional weight to people I love is understandable, but that preference does not automatically settle what is fair to everyone involved.
I do not think I need to abandon my concern for consequences to take those objections seriously.
I do need to recognize that my first way of evaluating a situation may leave something out.

Our small-group discussion of care ethics was particularly enlightening for this reason.
We debated whether companies could care, and spent some time sorting out what we meant by the words themselves.
We emerged with a greater understanding, and I enjoyed having that sort of intellectual exchange about the readings.
The discussion gave me a useful distinction to think about.
A company does not have feelings in the same way a person does, but its practices can still be evaluated according to whether they notice needs, assign responsibility, provide competent help, and respond to the people affected.
I do not think the semantic question is pointless, but it matters what follows from the definition.
Calling an organization caring is not worth very much if there is no way for the people receiving its care to say that it has misunderstood them.

\subsection{Assessment Reflection and Active Experimentation}
My 4GS self-ratings were 15 in giving, 12 in growth, 15 in grit, 14 in gratitude, and 17 in practice.
I asked my roommates and close friends Maggie and Alex to evaluate me on the same scales, and although they didn't use the exact scorecard, I asked for them to rank those qualities.
Maggie gave me an 18, 14, 15, 17, and 17, respectively.
Alex gave me a 17, 12, 15, 17, and 19, respectively.
The clearest agreement was in grit, where all three of us rated me the same.
The most useful difference was in giving, where both of my friends seemed to think I was more giving than I considered myself, which identifies a strength I wasn't entirely confident in.

To translate the two selected values from earlier, trust and temperance, into practice, I tried reminding myself to sit with my thoughts and evaluate actions before executing as many times as possible in a day.
The intended effects were fewer moments I would regret or look back on with the thought, ``why did I do that?''
What happened was some nicely thought out moments, where I did appreciate having taken the extra time to think, but mostly a lot of overthinking.
I found myself actually feeling more anxious than usual, and worrying about small and minute problems that didn't necessarily impact my life on an important level.

I adjusted the experiment by stopping the requirement that I pause deliberately in every conversation.
In fact, pausing before responding is something my father taught me as a child, so the habit itself was not new.
What changed was that I had turned it into a task I was trying to perform continuously.
The result suggests that more conscious monitoring is not always the same as better judgment.
For now, I want to use the pause particularly when I can feel that an emotional reaction is likely to determine what I say, rather than treating every ordinary interaction as a potential mistake.
I have changed the rule, but I cannot yet claim that this has solved the problem.

\subsection{Closing Ethical Reflection}
The quality I most want to develop is still temperance, though I mean something more specific by that than simply being quiet or appearing controlled.
I would like the things that happen to me, good or bad, to have less power to pull me into large emotional swings.
Those swings make it difficult to remain level-headed.
At the same time, my initial response to my father does not need to become another reason to judge myself harshly.
There is a difference between having a painful reaction and choosing to respond in a way that harms him.
The work is in how I handle that reaction, and in learning to evaluate his actions without requiring either him or myself to be morally perfect.

\section{Week 5. Practicing Listening \textit{(24 September 2026)}}

\subsection{Concrete Experience}
This summer, my friend group imploded.
My roommate since freshman year and closest friend, A, said something extremely unkind to another one of our friends, C, and his subsequent handling of the situation only served to amplify the issue.
Everyone but me was in Cleveland during this fallout, while I was doing research at Carnegie Mellon University in Pittsburgh, so I only ever got to hear about the issue over Google Meet.
When I came back for this school year, I was harboring a lot of anger and resentment against A that hadn't had a chance to be voiced or processed.
A few weeks into the semester, A and I were talking and he told me he'd been feeling underappreciated and ``like a pest.''
He told me, for the first time since it happened, the whole story from the summer from his perspective; he told me he still felt awful about what happened and his lackluster apology; he told me he was trying extremely hard to do better, and that in fact he and C were on much better terms now.
It turns out that S---another friend of ours---and I were actually the only ones who were still actively angry with him, despite being less directly involved than the others, and that A could feel our distrust of him despite us never voicing it.

I wanted to be upset with him, to tell him he wasn't allowed to be the victim in a situation he had created, that what he had done was unacceptable and warranted the harsh response from the rest of us.
But it was clear that what A needed in the moment was support from a friend, and that my anger was mostly unfair.
I felt the urge multiple times to interrupt him, to shout at him, but I knew that wasn't fair.

\subsection{Reflective Observation}
What strikes me in retrospect is that I was preparing to answer a version of the situation that was no longer entirely current.
I knew what had happened over the summer, but I had not been present for all of the conversations afterward.
I was still holding on to my anger while A and C had already begun to address the conflict between them.
That does not mean my original reaction was meaningless, or that what A did was acceptable.
It does mean that my distance from the situation had not made my judgment more objective.

I offered A my unwavering support as a friend and reminded him that I loved him.
I told him that I had heard about what happened over the summer, and admitted that I had initially felt resentful of him.
I also told him that I had come to see much of this as a boiling pot of unaddressed emotions that was not helping anyone.
If the other friends were willing to move toward reconciliation, I was too.
It was not fair to him for me to keep harboring that anger without ever dealing with it.
Our relationship now feels closer to what it used to be, warmer and less strained.

I think the important change was that I stopped treating support for him as necessarily incompatible with disapproval of what he had done.
He was not asking me to declare his conduct harmless.
He was telling me that he felt awful, that he was trying to improve, and that our distrust was affecting him.
I could acknowledge those things without deciding on C's behalf whether forgiveness was owed.
That distinction made it possible to respond to the friend in front of me rather than continuing an argument I had been carrying around in my head.

\subsection{Abstract Conceptualization and Reading Reflection}
Nichols \parencite{nichols2009listening} explains listening as the act of temporarily setting aside one's own agenda to understand another person's thoughts and feelings.
The difficulty can often be an internal demand to do something; to correct a detail, solve the problem, reassure the speaker, defend oneself, or relate by telling a personal story.
These responses come from a place of good intentions, but they redirect the conversation before the speaker has had a chance to communicate what matters.

The reading focuses on making listening active.
Encouraging elaboration and checking interpretation require attention and judgment.
Silence alone can leave a speaker uncertain as to whether their words are falling on deaf ears, but too much active listening can bulldoze their words.
This means that listening is possible even in disagreement, like my conversation with A.
I can make sure I accurately understand a person's experience before deciding to evaluate their conclusion.
An important point Nichols makes is that this doesn't require infinite emotional availability or the abandonment of my own personal boundaries.

Another useful part of Nichols's account is that conversations can develop patterns which neither person intends.
If a speaker senses that the listener is withdrawing or waiting for them to finish, they may become more insistent, which in turn can make the listener more impatient.
Seen that way, the difficulty is not always a defect located entirely in one person.
It can be a response to the relationship being created in the conversation.
Asking a question or checking what I have understood can interrupt that pattern, provided I am actually willing to hear an answer that changes my interpretation.
Otherwise, even a technically open question can become a way of steering the conversation toward what I already believe.

Connecting these ideas to care ethics, Engster and Hamington \parencite{engsterHamington2015introduction} emphasize responsiveness to the person receiving care.
If the carer decides what the cared for needs before hearing them, the help may serve the carer's desire to feel useful more than the circumstances of the cared for.
Listening gives the other person authority over the description of their experience.
In my conversation with A, the practical issue wasn't about whether I supplied good answers or not, but whether I understood the problem he was trying to communicate.

I mostly agreed with the Nichols reading.
As a speaker, the guidelines Nichols laid out would in fact make me feel more understood, and the ideas seem quite tractable as a listener.
There are occasions where one's input is necessary, however, and I would not want to interpret this advice as a requirement to remain passive.
Especially in the case where the speaker is saying things which could be harmful to themselves or others, I think it is important to redirect or offer advice.
The new insight I'd like to test is that the urge to help can sometimes interrupt the information needed to help well.
A limitation is that no technique guarantees trust.
For example, asking leading questions can stifle someone's ability to answer organically.
This is an idea I incorporate into my academic life.
Since I work as a tutor and a TA, it's important that I help my students arrive at answers on their own and I don't just lead them to the answer.
\subsection{Active Experimentation}
The two skills I selected for practice are allowing elaboration before responding and checking understanding before advising.
My conversation with A was an obvious chance to practice these, but I had further opportunity on a phone call with my mother a few days later.
She called to tell me about some stressful situations she'd been experiencing at school---she's a current PhD student and adjunct faculty member at Point Park University---and about the difficulties of taking care of her aging father, my grandfather.
When my mother and I usually talk, we're in the habit of finishing each other's sentences and flippantly offering advice.
I took this particular chance to practice allowing full elaboration of her problems without offering any input, then explicitly asking if she was looking for practical advice or just affirmation of her difficulties.
In fact, it turns out that in this particular case she wasn't looking for actionable ideas, she was looking for a thoughtful, listening ear.
Moreover, I believe I got to hear more details about how she was feeling, to more deeply empathize with her scenario, than I would have if I had simply treated her situation like I normally would.
I'm hoping to make this a habit, and I'll try to remind myself about this lesson during any conversations of importance.

\subsection{Closing Ethical Reflection}
The conversations with A and my mother clarified different parts of the same problem.
With A, I needed to make room for an account that complicated my anger.
With my mother, I needed to stop assuming that being useful meant supplying a solution.
In both cases, the other person had information about what they needed that I could not obtain by thinking harder on my own.
Listening therefore fits with my desire to be dependable, but it also asks something of me.
That is, I have to let the person depending on me help define what that dependability should look like.

\section{Week 7. Ethical Leadership Development \textit{(6 October 2026)}}

\subsection{Concrete Experience}
As I've mentioned before, I'm a TA for the introductory calculus classes.
Every Wednesday night I offer office hours for my students, and on Thursday mornings I lead a recitation and proctor a quiz (or exam) which I subsequently grade.
This year, I've been lecturing using the chalkboard, as opposed to projecting my writing on the document camera like I used to.
I made this decision for two reasons.
First of all, my recitation is in Strosacker auditorium, a massive venue, which happens to have a very large chalkboard.
Second of all, I wanted practice lecturing on a chalkboard, since I'm hoping to have a future in academia, and particularly in teaching.
I recently asked my students to fill out a feedback form regarding my instruction, in particular their preference of chalkboard vs projected notes, and was surprised to find a slight majority preferred the latter.

\subsection{Reflective Observation}
Despite the fact that I had personal reasons for preferring a chalkboard lecture, I can't deny the fact that my students prefer the projection of notes.
An ethical leader listens to input, adapts to the situation, and improves their leadership as such.
I've started lecturing by projecting my notes, and a few of my students have thanked me after class for making the change.
This is a clear example of how swallowing one's pride and listening attentively to those who depend on them can immediately improve the situation for both the leader and those being led.

\subsection{Abstract Conceptualization and Reading Reflection}
The Case and Smith definition of integrity \parencite{caseSmith2013integrity} connects words and actions, stated and enacted values, and consistency across situations and time.
It also involves ethical principles, transparency, and accepting consequences under pressure.
This makes integrity more demanding than doing whatever one has previously said.
The content of the commitment matters, and so does the willingness to let others assess whether it has been honored.

Newman \parencite[ch. 1]{Newman2019building} helps explain why this standard can be difficult to enact.
Recognizing an ethical issue, judging what should be done, prioritizing it over competing incentives, and carrying the decision through are distinguishable tasks.
I might understand a responsibility, but postpone action because acknowledgment would threaten my image or create an immediate inconvenience.
Locating the difficulty makes development more specific.
The needed change might be earlier recognition, clearer reasoning, greater motivation, or a practiced way to speak and act.

Together, these readings frame growth as a relationship between judgment, conduct, and context.
Revising an unrealistic promise openly can protect integrity better than silently preserving its appearance.
Likewise, revising a belief after uncovering evidence can be responsible, provided the reason for the change can be justified.
The question is whether the revision improves the ethical quality and reliability of conduct, or merely avoids an unwelcome cost.

My reaction to these accounts was almost complete agreement, especially as it related to my situation as a TA.
An insight that they offer is that trust depends partly on making limitations visible early enough for others to respond, like how I offered my students a chance to get the instruction that best suited them early enough in the semester to make a difference.

\subsection{Assessment Reflection}
From early in the semester, my values assessment identified my top five values, in order, as friendship, intellectualism, health, competency, and involvement.
I see these in my humanistic approach to instruction, although I struggle to enact them consistently in some professional scenarios.
The process of choosing them was useful partly because I could not simply keep everything I liked.
I had initially selected sixteen values, and narrowing the list forced me to distinguish a general priority from a particular expression of that priority.
Health was, and remains, one of the more aspirational choices.
I value it, but it is often the first thing I sacrifice under stress.

My Philosophical Orientation Questionnaire results were particularly consistent with this ranking.
I scored zero on the pragmatic lens, twenty-six on the intellectual lens, and thirty-four on the human lens, corresponding to the first, ninetieth, and ninety-fourth percentiles on the questionnaire's chart.
Friendship and intellectualism line up fairly directly with those last two results.
I still prefer examining the intellectual and interpersonal dimensions of a choice to declaring my first instinct common sense.
That being said, the pausing experiment gave me a concrete reason not to dismiss the value of execution.
An analysis which leaves me more anxious and no better able to act has not necessarily made the decision more ethical.

The VIA results require more qualification.
My highest-ranked strengths were judgment, kindness, love of learning, social intelligence, teamwork, humor, and appreciation of beauty and excellence.
The lowest were zest, hope, perseverance, self-regulation, love, and spirituality.
I remain skeptical of taking a quick self-report survey as a complete account of character, particularly when it seems to conflate a difficulty expressing something with its absence.
The low placement of love, for example, does not fit how I understand my relationships.
Likewise, working hard to regulate myself is different from having no capacity for self-regulation.
The questionnaire is useful as a reason to ask questions, but I would not treat its ranking as a verdict.

The comparison with my 4GS feedback makes this particularly clear.
My friends and I agreed on my grit score, although perseverance ranked relatively low in the VIA results.
Those are different instruments and should not be treated as directly comparable measurements, but the contrast suggests a question worth asking: in which circumstances do I persist, and in which do discomfort or discouragement stop me?
The same applies to kindness and giving.
Both friends rated my giving more highly than I did, which is evidence that my own critical self-perception is not the only account worth considering.

\subsection{Active Experimentation}
The teaching change supplies a concrete application of this reflection.
I asked the students for feedback, found that a slight majority preferred projected notes, and changed the way I delivered the recitation.
A few students subsequently thanked me for making the change.
That is a modest result, but it is an actual result; the feedback altered what I did, and some of the people affected noticed.
It would be too much to claim that the change has improved their understanding or their grades, since I have not established either of those things.
What I can say is that my preferred way of practicing teaching was not necessarily the format they preferred for learning.

For future recitations, I want to keep asking whether the format is serving the students rather than assuming that one piece of feedback settles the question permanently.
My interest in developing chalkboard teaching skills has not disappeared, but I need to find ways of developing those skills that remain responsive to the audience.
The environment helped in this case.
I had the discretion to change formats, a way for students to express a preference, and an alternative that I already knew how to use.
Those conditions made acting on the feedback relatively straightforward.

\subsection{Closing Ethical Reflection}
I do not think changing a lecture format makes me an unusually ethical person.
I do think it is a useful example of the kind of leadership I am already practicing, even when I do not describe myself as wanting to be a leader.
Students depend on my preparation and judgment, and that creates a responsibility to examine whether my preferences are serving them.
Integrity here means remaining consistent in the purpose of helping them learn while being willing to change the method.
That is a form of consistency I find much more convincing than refusing to revise an earlier decision simply because it was mine.

\newpage
\section{Synthesis of Ethical Leadership Learning}

At the beginning of this journal, I had trouble choosing someone whose character I wanted to resemble.
Looking back, I think the difficulty was partly in the question I thought I was being asked.
I was looking for a person whose character I could endorse entirely, when what I actually had were people whose particular qualities I admired and whose mistakes I could also recognize.
My father is the clearest example.
I still admire his intelligence, empathy, and integrity, but his decision to keep his job loss from me does not fit neatly into the image I had of him.
The useful change in my thinking is that I do not need to resolve that conflict by deciding that either my admiration or my disappointment was mistaken.
Both tell me something about the relationship and about the conduct I want to practice myself.

Wood, Meister, and Liu's distinction between effectiveness and morality helped make that separation possible \parencite{wood2021goodbadevil}.
An accomplishment is not sufficient evidence of moral character, just as a professional setback does not establish a lack of it.
That matters to me because I place so much value on intelligence and competency.
My first response to my father's situation included disbelief that someone I considered so smart could have lost his job.
In retrospect, that reaction connected his abilities, his professional circumstances, and his character more tightly than the evidence justified.
I can question his decision to conceal the situation without treating the job loss itself as a moral failure.
Newman's distinction between competence, commitment, and character reinforces this point.
The qualities interact, but none is a substitute for the others \parencite[ch. 1]{Newman2019building}.

I see a similar problem in my response to A.
I was holding on to an account of what had happened over the summer while the people most directly involved were already moving toward reconciliation.
I do not think my anger was inexplicable, but I had allowed it to remain unexamined.
Once A told me how he had been feeling, I could recognize that maintaining my distance was itself affecting someone I loved.
Supporting him did not require me to approve of what he had done.
It required me to respond to his present effort and to acknowledge my own part in the strain between us.
Our relationship feels warmer now, which is more meaningful evidence of the conversation's value than simply saying that I became a better listener.

The care ethics reading and Nichols's account of listening gave me language for this distinction \parencite{engsterHamington2015introduction,nichols2009listening}.
If I decide in advance what someone needs, I can mistake my preferred response for care.
With A, that response would have been to repeat my disapproval; with my mother, it would have been to offer advice before asking whether she wanted it.
In both cases, listening allowed information into the conversation that my own reasoning could not supply.
Our small-group debate about whether companies could care extended this question beyond personal relationships.
I enjoyed the intellectual exchange, but what remains useful is the practical question underneath the terminology; can the people receiving care influence how their needs are understood?
That is also a question I can ask of myself as a friend or a TA.

My assessment results make these connections fairly unsurprising.
Friendship and intellectualism were my two highest values, and the human and intellectual lenses were by far my strongest POQ results.
I want to understand things carefully and to be useful to the people around me.
The problem is that either tendency can become less helpful when I assume that more of it is always better.
The deliberate pausing experiment is a good example.
I was trying to practice trust and temperance, but monitoring every conversation mostly made me overthink.
I have stopped imposing that requirement on every interaction.
The next step is to use the pause where it serves a purpose, particularly when a strong emotional response might determine what I say, rather than making ordinary conversation another test I might fail.

My Week 2 reflection makes this tendency toward perfection more concrete.
I still regret how I performed academically in high school, and I think some of my current determination comes from not wanting to repeat that experience.
There is something useful in that motivation, but avoiding every situation in which I might fail would also mean avoiding opportunities I actually want.
This semester's schedule is an experiment in accepting that risk.
I chose nineteen credits, including nine graduate-level credits, alongside research, two jobs, and club commitments because I was excited about the work.
So far, I feel good about applying my abilities, but I am also stressed and frequently overwhelmed.
I do not think either observation cancels the other, and I am not yet sure that the schedule was the right choice.

This connects the failure reading to both my assessment results and the pausing experiment.
Rea, Stoller, and Kolp's emphasis on character, competence, and support suggests that taking on more difficulty does not automatically produce more growth \parencite[ch. 1]{Rea2022better}.
Similarly, making myself deliberate over every conversation did not automatically produce better judgment.
In both cases, the useful question is whether the extra effort is helping me do something I value, and what it costs to sustain it.
Health is one of my stated priorities, but I have already noticed that it is often the first one I sacrifice under pressure.
That gives me a reason to evaluate the semester by more than whether I maintain a perfect academic record.
The experiment remains unfinished; feeling capable and interested in the work is encouraging, but simply not having crashed and burned is not sufficient evidence that the balance is sustainable.

Temperance is still my most important development need.
I would like both good and bad experiences to have less power to produce the large emotional swings that make it difficult for me to stay level-headed.
That is not the same as wanting to become indifferent to people or to what happens around me.
My relationships matter precisely because I care about them.
The distinction I need to practice is between taking a feeling seriously and allowing its intensity to settle what I should do.
The Rea, Stoller, and Kolp discussion of virtues as practiced capacities is useful here; knowing that I value temperance does not mean that I have already developed a reliable way to exercise it \parencite[ch. 2]{Rea2022better}.

The assessments also remind me to be careful about turning a difficulty into a complete account of myself.
My VIA results placed perseverance and self-regulation relatively low, while my friends' 4GS feedback confirmed my grit and rated my giving more highly than I did.
I do not think these differences can be resolved by declaring one instrument correct.
They make more sense as questions about context and about the gap between how much effort something takes internally and what another person experiences externally.
I can recognize that emotional regulation is difficult for me while also recognizing the effort and restraint I already exercise.
Likewise, the fact that health appears in my top five values does not mean that I consistently protect it when I am stressed.
That is a gap between valuation and execution worth continuing to work on.

My teaching provides a smaller but clearer example of translating a value into action.
I preferred the chalkboard partly because I wanted to develop a skill I expect to use in academia.
When students expressed a preference for projected notes, I changed the format, and some thanked me afterward.
The integrity connection is that the underlying commitment remained the same while the method changed \parencite{caseSmith2013integrity}.
If my purpose is to help students learn, then being consistent cannot mean insisting on a method regardless of their experience.
I want to continue asking for feedback while being precise about what it establishes.
A preference and a positive response are useful evidence, but they do not by themselves demonstrate improved learning.

I still do not particularly want to be the Guy.
I would rather be someone the Guy, my friends, or my students can depend on.
The journal has made it harder to treat that preference as a reason leadership character is irrelevant to me.
Dependability involves judgment about what to do when someone is upset, when I disagree, or when my preferred approach is not serving the people relying on me.
Writing these entries helped me notice a repeated tendency to judge through an idealized picture, whether of my father, a friend, or my own conduct.
My aim going forward is to keep the standards without requiring perfection---listen before settling the interpretation, pause when the emotion needs space, and revise a method when the response gives me reason to do so.
I cannot yet claim consistency in all of that, but it is a more concrete direction than simply wanting to become a better person.

\printbibliography[heading=bibintoc]
\end{document}
